\documentclass[11pt]{article}

\usepackage[margin=1in]{geometry}
\usepackage{amsmath, amssymb, bm}
\usepackage{enumitem}
\usepackage{hyperref}
\usepackage{lmodern}
\usepackage[T1]{fontenc}
\usepackage[utf8]{inputenc}

\hypersetup{
    colorlinks=true,
    linkcolor=blue,
    urlcolor=blue
}

\setlength{\parindent}{0pt}
\setlength{\parskip}{6pt}

\title{Mathematical Model Specification\\
\large Stochastic Control Routines for Optimal Execution and Market Making}
\author{}
\date{}

\begin{document}
\maketitle

This document states the two implemented stochastic-control models without relying on external notes. The project is synthetic and model-based; it is not a calibrated trading strategy or a production market-making system.

\section{Closed-Form Optimal Liquidation}

The liquidation state variables are:
\begin{itemize}[itemsep=2pt]
    \item $Q_t$: remaining inventory;
    \item $S_t$: unaffected midprice;
    \item $X_t$: cash account.
\end{itemize}

The control is the selling rate $a_t$, where positive values represent sales. Admissible selling rates are progressively measurable, nonnegative, and square-integrable over the horizon. For the nonnegative initial inventories used in this project, the closed-form optimizer below is nonnegative.

The parameters are terminal time $T$, midprice volatility $\sigma$, temporary impact coefficient $\kappa$, and terminal inventory penalty $\theta$.

The model uses no permanent market impact. The controlled dynamics are
\begin{align}
dQ_t &= -a_t\,dt, \\
dS_t &= \sigma\,dW_t, \\
\widehat{S}_t &= S_t-\kappa a_t, \\
dX_t &= a_t\widehat{S}_t\,dt
     = a_t(S_t-\kappa a_t)\,dt.
\end{align}

The terminal objective is
\begin{equation}
\sup_a \mathbb{E}\left[X_T + Q_T S_T - \theta Q_T^2\right].
\end{equation}

Equivalently, because cash is generated by the execution equation,
\begin{equation}
\sup_a \mathbb{E}\left[
\int_0^T a_t(S_t-\kappa a_t)\,dt
+ Q_T S_T - \theta Q_T^2
\right].
\end{equation}

Using the value-function form
\begin{equation}
V(t,S,x,q)=x+Sq+\gamma(t)q^2,
\end{equation}
the no-permanent-impact HJB reduces to the Riccati equation
\begin{equation}
\gamma'(t)=-\frac{\gamma(t)^2}{\kappa},
\qquad
\gamma(T)=-\theta.
\end{equation}

The solution is
\begin{equation}
\gamma(t)
=-\left(\frac{1}{\theta}+\frac{T-t}{\kappa}\right)^{-1}.
\end{equation}

The optimal selling rate is
\begin{equation}
a^*(t,q)
=-\frac{\gamma(t)}{\kappa}q.
\end{equation}

The resulting deterministic inventory path is
\begin{equation}
Q(t)
=Q(0)\frac{\kappa/\theta+T-t}{\kappa/\theta+T}.
\end{equation}

The finite terminal penalty does not impose the hard constraint $Q_T=0$. Residual terminal inventory remains unless the terminal penalty is taken to a limiting hard-liquidation regime.

\section{Symmetric Limit-Order Market Making}

The market-making state variables are:
\begin{itemize}[itemsep=2pt]
    \item $S_t$: midprice;
    \item $X_t$: cash account;
    \item $Q_t$: inventory.
\end{itemize}

The controls are quote distances $\delta_b$ and $\delta_a$, respectively the bid distance below the midprice and ask distance above the midprice. Admissible quote distances are predictable and nonnegative. Quotes that would move inventory outside the finite inventory grid are suppressed by setting the corresponding distance to infinity, which gives zero fill intensity.

The parameters are order size $\Delta$, baseline fill-arrival intensity $\lambda$, exponential fill-decay coefficient $\kappa$, midprice volatility $\sigma$, terminal inventory penalty $\alpha$, running inventory penalty $\phi$, and finite inventory bounds $q_{\min}$ and $q_{\max}$.

The bid and ask execution prices are $S_t-\delta_b$ and $S_t+\delta_a$. Controlled fill intensities are
\begin{align}
\lambda_b(\delta_b) &= \lambda e^{-\kappa\delta_b}, \\
\lambda_a(\delta_a) &= \lambda e^{-\kappa\delta_a}.
\end{align}

Let $N_t^b$ and $N_t^a$ denote the bid-fill and ask-fill counting processes. The controlled state dynamics are
\begin{align}
dS_t &= \sigma\,dW_t, \\
dQ_t &= \Delta\,dN_t^b-\Delta\,dN_t^a, \\
dX_t &= -\Delta(S_t-\delta_b)\,dN_t^b
      + \Delta(S_t+\delta_a)\,dN_t^a.
\end{align}

The objective is
\begin{equation}
\sup_{\delta_b,\delta_a}
\mathbb{E}\left[
X_T + Q_T S_T
- \alpha Q_T^2
- \int_0^T \phi Q_t^2\,dt
\right].
\end{equation}

The terminal condition is
\begin{equation}
v(T,S,x,q)=x+Sq-\alpha q^2.
\end{equation}

For the symmetric case implemented here, the Bellman equation is
\begin{align}
0
&= \partial_t v + \frac{1}{2}\sigma^2\partial_{SS}v - \phi q^2 \notag\\
&\quad + \lambda\sup_{\delta_b \ge 0}
e^{-\kappa\delta_b}
\left[
v(t,S,x-\Delta(S-\delta_b),q+\Delta)-v(t,S,x,q)
\right]\mathbf{1}_{q<q_{\max}} \notag\\
&\quad + \lambda\sup_{\delta_a \ge 0}
e^{-\kappa\delta_a}
\left[
v(t,S,x+\Delta(S+\delta_a),q-\Delta)-v(t,S,x,q)
\right]\mathbf{1}_{q>q_{\min}}.
\end{align}

Use the ansatz
\begin{equation}
v(t,S,x,q)=x+Sq+\theta(t,q).
\end{equation}

Because this ansatz is linear in the midprice $S$, the second derivative of $v$ with respect to $S$ is zero. Therefore $\sigma$ affects simulated marked-to-market outcomes, but it does not enter the transformed ODE or the quote policy in this symmetric model.

The bid and ask jump differences become
\begin{align}
v(t,S,x-\Delta(S-\delta_b),q+\Delta)-v(t,S,x,q)
&= \Delta\delta_b+\theta(t,q+\Delta)-\theta(t,q), \\
v(t,S,x+\Delta(S+\delta_a),q-\Delta)-v(t,S,x,q)
&= \Delta\delta_a+\theta(t,q-\Delta)-\theta(t,q).
\end{align}

Define
\begin{align}
g_b(t,q)
&= \frac{\theta(t,q+\Delta)-\theta(t,q)}{\Delta}, \\
g_a(t,q)
&= \frac{\theta(t,q-\Delta)-\theta(t,q)}{\Delta}.
\end{align}

In the interior-positive regime, where the unconstrained maximizers have positive quote distances, the first-order conditions give
\begin{align}
\delta_{b}^{*}(t,q) &= \frac{1}{\kappa}-g_b(t,q), \\
\delta_{a}^{*}(t,q) &= \frac{1}{\kappa}-g_a(t,q).
\end{align}

Equivalently,
\begin{align}
\delta_{b}^{*}(t,q)
&= \frac{1}{\kappa}
   - \frac{\theta(t,q+\Delta)-\theta(t,q)}{\Delta}, \\
\delta_{a}^{*}(t,q)
&= \frac{1}{\kappa}
   - \frac{\theta(t,q-\Delta)-\theta(t,q)}{\Delta}.
\end{align}

At the boundaries, the implementation suppresses quotes that would move inventory outside the grid:
\begin{equation}
\delta_{b}^{*}(t,q_{\max})=+\infty,
\qquad
\delta_{a}^{*}(t,q_{\min})=+\infty.
\end{equation}

The code checks that all finite quote distances are strictly positive for the chosen parameters. If the unconstrained first-order condition produced a nonpositive distance, the constrained-control problem would need to be solved with the nonnegativity constraint active; that case is outside this toy implementation.

After substituting the optimal quotes, the equation for $\theta$ is nonlinear. The implemented transformation is
\begin{equation}
\theta(t,q)=\frac{\Delta}{\kappa}\log W(t,q).
\end{equation}

This gives the linear finite-grid ODE
\begin{equation}
W'(t)+AW(t)=0,
\qquad
W(t)=\exp((T-t)A)W(T).
\end{equation}

The terminal condition and matrix entries are
\begin{align}
W(T,q) &= \exp\left(-\frac{\kappa}{\Delta}\alpha q^2\right), \\
A_{q,q} &= -\frac{\kappa}{\Delta}\phi q^2, \\
A_{q,q+\Delta} &= \lambda e^{-1}, \\
A_{q,q-\Delta} &= \lambda e^{-1}.
\end{align}

Missing off-diagonal terms at the inventory boundaries are omitted.

\section{Simulation Discretization}

The simulation approximates continuous-time Poisson fills over a time step $\Delta t$ with a first-event discretization. Given quote-dependent intensities
\begin{align}
\lambda_b &= \lambda e^{-\kappa\delta_b}, \\
\lambda_a &= \lambda e^{-\kappa\delta_a}, \\
\Lambda &= \lambda_b+\lambda_a,
\end{align}
the probability of at least one fill in the step is
\begin{equation}
p_{\mathrm{fill}}=1-\exp(-\Lambda\,\Delta t).
\end{equation}

Conditional on a fill, the side is sampled by relative intensity:
\begin{equation}
\mathbb{P}(\mathrm{bid}\mid\mathrm{fill})=\frac{\lambda_b}{\Lambda},
\qquad
\mathbb{P}(\mathrm{ask}\mid\mathrm{fill})=\frac{\lambda_a}{\Lambda}.
\end{equation}

If a boundary quote is suppressed, its fill intensity is set to zero. If $\Lambda=0$, no fill is sampled. This discretization allows at most one fill per time step.

When a bid fill occurs,
\begin{align}
X_{n+1} &= X_n-\Delta(S_n-\delta_b), \\
Q_{n+1} &= Q_n+\Delta.
\end{align}

When an ask fill occurs,
\begin{align}
X_{n+1} &= X_n+\Delta(S_n+\delta_a), \\
Q_{n+1} &= Q_n-\Delta.
\end{align}

The pathwise running inventory penalty is approximated by the left-point sum
\begin{equation}
R_T^{\Delta t}
= \sum_{n=0}^{N-1}\phi Q_n^2\,\Delta t.
\end{equation}

The simulator reports both the terminal marked-to-market diagnostic
\begin{equation}
X_N+Q_N S_N-\alpha Q_N^2
\end{equation}
and the objective-consistent discretized payoff
\begin{equation}
X_N+Q_N S_N-\alpha Q_N^2-R_T^{\Delta t}.
\end{equation}

\end{document}
